\chapter{Линейные преобразования для получения решений пространственной задачи встречи с учётом эллиптичности орбит}

% =========================================================
% Рабочий каркас главы 3.
% Это НЕ финальная редакция главы, а удобная структура для
% аккуратного переноса материалов из ВКР и презентации.
%
% Принцип заполнения:
% - сначала переносится опорный текст, формулы и иллюстрации из ВКР;
% - затем материал перестраивается по логике кандидатской диссертации;
% - служебные комментарии и указания на источники удаляются только
%   на позднем этапе шлифовки.
% =========================================================

% Общая опора:
% ВКР: глава 4, с. 40--59
% Презентация: слайды 26--32

\section{Проблема перехода от плоской задачи к пространственной постановке}
% Опора:
% ВКР: глава 4, с. 40--59
% Презентация: слайд 26
%
% Здесь ввести мотивацию главы:
% - результаты главы 2 относятся к плоской круговой модели;
% - для баллистического проектирования нужны решения для
%   пространственных перелётов между эллиптическими орбитами;
% - прямой пересчёт таких решений затруднён;
% - требуется специальный класс преобразований, сохраняющий
%   полезную структуру начального приближения.
%
% Желательно явно сформулировать:
% - в чём состоит "проблема деформации";
% - почему одного поворота системы координат недостаточно.
%
% Черновой текст:

\section{$R$-преобразование как специальный класс линейных преобразований}
% Опора:
% ВКР: §4.1--4.2, с. 40--51
% Презентация: слайд 27
%
% Здесь нужно ввести:
% - матрицу линейного преобразования R;
% - её специальную структуру;
% - смысл преобразования для фазовых и сопряжённых переменных;
% - сценарии использования с изменением и без изменения фазовых
%   переменных.
%
% Важно:
% - не ограничиваться записью формул;
% - сразу связать R-преобразование с практической задачей перехода
%   от аппроксимированного многообразия к более реалистическим
%   граничным условиям.
%
% Черновой текст:

\section{Первые интегралы и ограничения на линейные преобразования}
% Опора:
% ВКР: §4.1, с. 40--45
% Презентация: слайды 26--27 как контекст
%
% Здесь собрать:
% - какие первые интегралы и структурные свойства важно сохранить;
% - какие ограничения это накладывает на допустимые преобразования;
% - почему специальный вид матрицы R выделяется среди общего класса
%   линейных преобразований.
%
% Этот раздел можно писать более строго и математично, чем
% презентационный материал, но без перегрузки второстепенными деталями.
%
% Черновой текст:

\section{Пространственное $R$-преобразование}
% Опора:
% ВКР: §4.3, с. 52--54
% Презентация: слайд 28
%
% Здесь нужно последовательно показать:
% - как задаются целевые наклонение i и долгота восходящего узла \Omega;
% - как по конечным формулам восстанавливается матрица R;
% - как преобразуются параметры начального приближения;
% - в чём состоит полезный эффект пространственного преобразования.
%
% Обязательно оставить:
% - замечание о том, что при таком преобразовании меняются
%   радиус конечной орбиты и угловая дальность;
% - визуальный пример соответствующей деформации.
%
% Черновой текст:

\section{Плоское $R$-преобразование}
% Опора:
% ВКР: §4.4, с. 55--59
% Презентация: слайд 29
%
% Здесь собрать:
% - постановку задачи для эксцентриситета e и аргумента перицентра \omega;
% - получение матрицы R по конечным формулам;
% - интерпретацию этого преобразования;
% - возникающую деформацию большой полуоси и угловой дальности.
%
% Важно:
% - текстом явно развести пространственное и плоское преобразования,
%   чтобы они не сливались в один технический блок.
%
% Черновой текст:

\section{Пересчёт решения на новую дату старта при эллиптической начальной орбите}
% Опора:
% ВКР: глава 4, с. 40--59
% Презентация: слайд 30
%
% Здесь нужно объяснить:
% - почему эллиптичность начальной орбиты создаёт отдельную проблему;
% - почему простое вращение системы координат даёт большую невязку;
% - чем R-преобразование лучше обычного поворота;
% - как это проявляется в численных экспериментах.
%
% Желательно явно сохранить тезис из презентации:
% - чувствительность под воздействием R-преобразования оказывается
%   на несколько порядков меньше, чем под воздействием матрицы поворота.
%
% Черновой текст:

\section{Проблема деформации траекторий и её компенсация}
% Опора:
% ВКР: глава 4, с. 40--59
% Презентация: слайд 31
%
% Здесь центральный смысл главы:
% - преобразования позволяют контролировать новые параметры орбиты,
%   но деформируют радиус / полуось и угловую дальность;
% - эту проблему нельзя игнорировать;
% - она устраняется за счёт возврата к аппроксимированному
%   многообразию главы 2 и подбора исходной точки на нём.
%
% Здесь удобно показать:
% - пример с изменённым значением \kappa;
% - идею компенсации деформации;
% - итоговый переход к контролю над всеми параметрами перелёта.
%
% Черновой текст:

\section{Конечные формулы для связи элементов матрицы $R$ и кеплеровых элементов конечной орбиты}
% Опора:
% ВКР: глава 4, с. 40--59
% Презентация: слайды 27--31
%
% Здесь собрать в одном месте:
% - какие именно конечные формулы получены методом символьной регрессии;
% - какие параметры они связывают;
% - какова их роль в алгоритме построения начального приближения.
%
% Этот раздел лучше держать компактным и инструментальным:
% подробные громоздкие выражения при необходимости можно увести
% в приложение или оставить как ссылку на приложение.
%
% Черновой текст:

\section{Роль линейных преобразований в общей схеме баллистического проектирования}
% Опора:
% ВКР: глава 4, с. 40--59
% Презентация: слайды 26--32
%
% Здесь важно подвести мост к главе 4:
% - линейные преобразования не являются самостоятельной целью;
% - они замыкают связку между аппроксимированным многообразием
%   и прикладным алгоритмом выбора начального приближения;
% - благодаря им результаты плоской модели становятся пригодными
%   для практического баллистического проектирования.
%
% Черновой текст:

\section*{Выводы к главе 3}
% Опора:
% ВКР: глава 4, с. 40--59
% Презентация: слайд 32
%
% Здесь потом кратко зафиксировать:
% 1. введение специального класса линейных преобразований -- R-преобразования;
% 2. получение конечных формул, связывающих элементы матрицы R
%    с кеплеровыми элементами конечной орбиты;
% 3. выделение пространственного и плоского сценариев преобразования;
% 4. обнаружение и интерпретацию деформации траекторий;
% 5. способ компенсации деформации за счёт использования
%    аппроксимированного многообразия главы 2.
%
% Черновой текст:
